\documentclass[conference, a4paper]{IEEEtran}
\usepackage{cuted}
\setlength{\stripsep}{6pt}
%\documentclass[conference, a4paper, 10pt]{ieeeconf}

% \overrideIEEEmargins
\IEEEoverridecommandlockouts
% The preceding line is only needed to identify funding in the first footnote. If that is unneeded, please comment it out.
\usepackage{cite}
\usepackage{amsmath,amssymb,amsfonts}
\usepackage{graphicx}
\usepackage{textcomp}
\usepackage{xcolor}
\usepackage{float}
\usepackage{placeins}
\usepackage{booktabs}
\usepackage{array}
\usepackage{subcaption}
\def\BibTeX{{\rm B\kern-.05em{\sc i\kern-.025em b}\kern-.08em
    T\kern-.1667em\lower.7ex\hbox{E}\kern-.125emX}}

% Figures live in the Figures/ folder
\graphicspath{{Figures/}}

\begin{document}

\title{Body-Motion Control of a Simulated Aerial Swarm under First-Person View
\thanks{This work was supported by the Swiss National Science Foundation under Grant 200020\_212077.}
\thanks{This work involved human subjects or animals in its research. Approval of all ethical and experimental procedures and protocols was granted by the EPFL Human Research Ethics Committee under Application No.\ HREC 092-2023.}
}

\author{\IEEEauthorblockN{Yang Chen, Darius Giannoli, and Dario Floreano}
\IEEEauthorblockA{\textit{Laboratory of Intelligent Systems} \\
\textit{\'Ecole Polytechnique F\'ed\'erale de Lausanne (EPFL)}\\
Switzerland\\
\texttt{chenyang@ai.iit.tsukuba.ac.jp},\\
\texttt{darius.giannoli@epfl.ch}, \texttt{dario.floreano@epfl.ch}}
}

\maketitle

\begin{abstract}
Aerial swarms require intuitive, continuous group-level commands under
immersive first-person view (FPV). Conventional radio-control transmitters
distribute commands across sticks and auxiliary controls, potentially
constraining simultaneous adjustment of swarm motion, view direction, and
formation size. We present a participant-calibrated upper-body interface for
five-degree-of-freedom FPV teleoperation of a simulated aerial swarm. The
interface maps torso inclination to planar velocity, mean hand height and
inter-hand separation to vertical velocity and inter-agent spacing, and head
yaw to FPV yaw rate. We compared the complete interface with a conventional
transmitter in a within-subject study in which 14 participants controlled a
simulated 15-agent swarm through three-dimensional gap-passing courses. The
body-motion interface reduced mean completion time by 19.7\% and centroid
travel distance by 6.6\%. No interface differences were detected in gate-centering accuracy, collection
yield, swarm survival, connectivity, total workload, or usability. Within the
tested simulation, these findings support an efficiency advantage for the
complete implemented body-motion interface and demonstrate its capacity for
concurrent multi-degree-of-freedom input. Evaluation under matched signal
processing and physical flight is required to isolate modality effects and
assess generalizability.
\end{abstract}

\begin{IEEEkeywords}
aerial swarm, body-machine interface, teleoperation, first-person view,
human--swarm interaction
\end{IEEEkeywords}

% =====================================================================
\section{Introduction}

Aerial swarms can provide parallel sensing, spatial coverage, and redundancy
for applications such as search and rescue and infrastructure inspection
\cite{ben_focaldrone,macchini2021personalized}. Even when decentralized
controllers handle local coordination and collision avoidance, complex or
safety-critical missions may still require a human operator to provide
high-level motion and formation commands. To preserve scalability, a single operator should command the collective
rather than control every agent individually
\cite{kolling2016survey,hocraffer2017meta}. The operator must therefore coordinate several swarm degrees of
freedom (DOF) at once---including planar translation, altitude, view direction,
and expansion or contraction---while perceiving the remote environment.

\begin{figure*}[t]
\centering
\includegraphics[width=\textwidth]{Figures/main.png}
\caption{Overview of this work. The same aerial swarm (center) is teleoperated
over five color-coded degrees of freedom (DOF) either with a Taranis
radio-control transmitter (left) or with the proposed upper-body interface
(right). In the body-motion interface, head yaw sets the view direction, torso
lean the planar motion, arm lift the vertical motion, and arm
flexion/extension the swarm's expansion/contraction. We compare the two
interfaces under first-person view on a 3D gap-passing course.}
\label{fig:overview}
\end{figure*}

The operator can perceive the swarm either from a third-person view (TPV), an
external overview of the group, or from a first-person view (FPV), the
egocentric image supplied by one swarm agent (Fig.~\ref{fig:views}). TPV makes
the swarm geometry directly visible, but obtaining a live external view in an
uninstrumented deployment may require additional sensing, reconstruction, or a
dedicated viewing platform. FPV can instead be obtained from an onboard camera,
but it limits global awareness and moves with the focal drone. Commands must
therefore be expressed in a changing egocentric frame, and the view direction
itself becomes a controlled DOF \cite{ben_focaldrone,macchini2021impact}. We
focus on FPV because it is compatible with onboard perception, while recognizing
that it places greater demands on the command interface.

Existing focal-drone FPV interfaces retain a conventional hand-held
transmitter, distributing planar and vertical motion and view yaw across two
thumbsticks while assigning swarm expansion to a separate knob
\cite{ben_focaldrone}. Although familiar, this arrangement may constrain the
simultaneous manipulation of all five DOF. The limitation is particularly
relevant when negotiating three-dimensional openings, where translation,
altitude, heading, and formation size may need to change concurrently. Thus,
even after an FPV viewpoint is established, the input device remains a
potential bottleneck.

Upper-body motion instead provides several continuous input dimensions. Head,
torso, and arm movements can be assigned to distinct swarm DOF, allowing
several commands to be expressed concurrently. Moreover, body lean can be
mapped to planar motion in the current focal-drone frame, preserving spatial
congruence as the view rotates or transfers between agents. Upper-body
interfaces have already enabled accurate and immersive control of
single drones \cite{rognon2018flyjacket,miehlbradt2018datadriven}, motivating
their extension to group-level control.

Human--swarm interaction reduces a high-dimensional collective to
operator-level commands, making interface design and workload central to
scalability \cite{kolling2016survey,hocraffer2017meta}. Gesture-based systems
have enabled users to manipulate meaningful swarm-shape DOF
\cite{alonsomora2015gesture,suresh2019gesture}. For aerial swarms,
personalized hand motion has provided continuous group-level control under TPV
\cite{macchini2021personalized}, SwarmTouch added tactile feedback
\cite{tsykunov2019swarmtouch}, and Krátký et al. used contactless gestures to
command a three-UAV formation \cite{kratky2025gesture}. These systems use
external viewpoints and/or discrete high-level gestures. Under FPV,
upper-body interfaces have focused on individual drones
\cite{rognon2018flyjacket,miehlbradt2018datadriven,macchini2024datadriven},
while swarm interfaces have retained conventional hand-held controls
\cite{ben_focaldrone}. Upper-body motion offers several continuous input
dimensions that can be mapped in the moving focal-drone frame. To the best of
our knowledge, no prior system combines continuous upper-body control of
translation, altitude, view direction, and formation size with an FPV from
within the swarm.

In this work, we design and evaluate an upper-body interface for FPV
teleoperation of an aerial swarm (Fig.~\ref{fig:overview}). We investigate
whether it enables faster task execution and smoother, more coordinated
multi-DOF commands than a conventional transmitter, and whether any gains occur
without a detectable degradation in task accuracy, swarm integrity, workload,
or usability. The contributions are threefold: (i) a spatially congruent
mapping from head, torso, and arm motion to five continuous swarm DOF under FPV;
(ii) a participant-calibrated implementation that remains aligned with the
current focal-drone frame; and (iii) a within-subject comparison with a
radio-control transmitter using task performance, control behavior, swarm
integrity, and subjective measures.

% =====================================================================
\section{Body-Motion Control Interface}
% --- 2a. Methodology: the interface and the body-to-swarm mapping ---
\subsection{System Overview}
\label{sec:system}

The implemented system couples an upper-body input layer to a collective
controller for a 15-agent aerial swarm simulated in Unity, as summarized in
Fig.~\ref{fig:overview}. The operator receives a first-person view (FPV) from
a focal drone selected automatically at the rear of the largest connected
component -- following the focal-drone framework evaluated
by Jarvis et al.~\cite{ben_focaldrone}. The operator commands the swarm as a
collective rather than controlling individual agents.

At each simulation frame, the input layer generates the five-dimensional
command vector
\begin{equation}
\mathbf{u}_t =
\begin{bmatrix}
v_{f,t}^{\mathrm{ref}} &
v_{l,t}^{\mathrm{ref}} &
v_{h,t}^{\mathrm{ref}} &
\dot{\psi}_{t}^{\mathrm{cmd}} &
\rho_t^{*}
\end{bmatrix}^{\mathsf{T}},
\label{eq:swarm_command}
\end{equation}
where $v_{f,t}^{\mathrm{ref}}$, $v_{l,t}^{\mathrm{ref}}$, and
$v_{h,t}^{\mathrm{ref}}$ are the forward, lateral, and vertical reference
velocities expressed in the current focal-drone frame;
$\dot{\psi}_{t}^{\mathrm{cmd}}$ is the commanded FPV yaw rate; and
$\rho_t^{*}$ is the target inter-agent spacing. The translational and yaw
components are signed rate commands, whereas $\rho_t^{*}$ is an absolute
formation-spacing setpoint.

A decentralized controller based on the Olfati--Saber flocking algorithm
\cite{olfatisaber2006flocking} converts these group-level commands into
agent motion while maintaining local cohesion, velocity alignment, and
obstacle avoidance. The human therefore specifies collective translation,
view direction, altitude, and formation size, while local coordination remains
autonomous. Both the body-motion and transmitter conditions drive this same
command vector and downstream swarm controller; they differ in their physical
input mappings and associated signal-processing pipelines.

In the body-motion condition, a Meta Quest~3S headset measures head yaw and
provides the FPV display while the two Quest controllers provide three-dimensional hand
positions. An LPMS-B2 inertial measurement unit (LP-Research Inc.) measures
torso inclination. The following subsections define the participant-calibrated
mapping from these measurements to Eq.~\eqref{eq:swarm_command}, followed by
the transmitter baseline.

\subsection{Body-to-Swarm Mapping}
\label{sec:mapping}

\begin{figure}[t]
\centering
\includegraphics[width=\columnwidth]{Figures/fig_mapping_technical_v2.png}
\caption{Mapping from measured upper-body motion to collective swarm
commands.}
\label{fig:mapping}
\end{figure}

At time $t$, the interface measures the sagittal and lateral torso
inclinations $(\theta_t,\phi_t)$, headset yaw $\psi_t$, mean vertical position
of the two controllers $h_t$, and their Euclidean separation $s_t$
(Fig.~\ref{fig:mapping}). Rate control was used for translation and view
rotation so that a bounded body posture could generate sustained motion.
Swarm spacing was instead mapped absolutely, providing a direct correspondence
between controller separation and formation size.

\emph{Rate-command normalization:}
Torso inclination, controller height, and head-yaw deviation were normalized
using the participant-specific limits obtained through the calibration in
Section~\ref{sec:calibration}. For a measurement $q_t$ with lower, neutral,
and upper calibration values $(q^{-},q^{0},q^{+})$ and deadzone half-width
$d$, let $q_{\mathrm{L}}=q^{0}-d$ and $q_{\mathrm{U}}=q^{0}+d$. The normalized
command was
\begin{equation}
\mathcal{D}(q_t)=
\begin{cases}
\displaystyle
\max\!\left(-1,\frac{q_t-q_{\mathrm{L}}}
{q_{\mathrm{L}}-q^{-}}\right),
& q_t<q_{\mathrm{L}},\\[7pt]
0,
& q_{\mathrm{L}}\leq q_t\leq q_{\mathrm{U}},\\[3pt]
\displaystyle
\min\!\left(1,\frac{q_t-q_{\mathrm{U}}}
{q^{+}-q_{\mathrm{U}}}\right),
& q_t>q_{\mathrm{U}}.
\end{cases}
\label{eq:deadzone_mapping}
\end{equation}
For subsequent expressions, the notation
$\mathcal{D}(q_t;q^{-},q^{0},q^{+},d)$ identifies the calibration bounds and
deadzone used in this mapping. Each side of the neutral posture was therefore
scaled independently to accommodate asymmetric ranges of motion.

\emph{Planar and vertical translation:}
Sagittal and lateral torso inclinations generated forward and lateral
commands, respectively, while mean controller height generated the vertical
command:
\begin{align}
c_{f,t} &=
\mathcal{D}(\theta_t;\theta^{-},0,\theta^{+},5^{\circ}),
\nonumber\\
c_{l,t} &=
\mathcal{D}(\phi_t;\phi^{-},0,\phi^{+},5^{\circ}),
\nonumber\\
c_{h,t} &=
\mathcal{D}(h_t;h_{\min},h_0,h_{\max},0.05~\mathrm{m}),
\label{eq:body_rate_commands}
\end{align}
where positive values command forward, rightward, and upward motion,
respectively.

Defining $\mathbf{c}_{p,t}=[c_{f,t},c_{l,t}]^{\mathsf T}$, the planar command
was low-pass filtered according to
\begin{equation}
\bar{\mathbf{c}}_{p,t}
=(1-\alpha_t)\bar{\mathbf{c}}_{p,t-1}
+\alpha_t\mathbf{c}_{p,t},
\qquad
\alpha_t=\min\{1,7\Delta t\},
\label{eq:torso_filter}
\end{equation}
where $\Delta t$ is the frame interval. The vertical command was filtered
separately using a critically damped filter with a 0.1-s smoothing time.
These filters attenuated IMU and controller-tracking noise before the
commands were supplied to the swarm controller.

Let $\mathbf{e}_{f,t}$, $\mathbf{e}_{r,t}$, and $\mathbf{e}_{u,t}$ denote the
forward, right, and up axes of the current focal drone. The reference velocity
applied to the swarm was
\begin{equation}
\mathbf{v}^{\mathrm{ref}}_t =
v_{\max}\left(
\bar c_{f,t}\mathbf{e}_{f,t}
+\bar c_{l,t}\mathbf{e}_{r,t}
+\bar c_{h,t}\mathbf{e}_{u,t}
\right),
\qquad
v_{\max}=10~\mathrm{m\,s^{-1}},
\label{eq:reference_velocity}
\end{equation}
where the bars denote the corresponding filtered commands. Consequently,
torso commands remained aligned with the displayed FPV heading, including
after the viewpoint transferred between focal drones.

\emph{Inter-agent spacing:}
Controller separation specified the target inter-agent spacing:
\begin{align}
e_t &=
\mathrm{clip}\!\left(
\frac{s_t-s_{\min}}{s_{\max}-s_{\min}},0,1\right),
\nonumber\\
\tilde{\rho}^{*}_t &=
\rho_{\min}+e_t(\rho_{\max}-\rho_{\min}),
\label{eq:spread_mapping}
\end{align}
where $\rho_{\min}=1~\mathrm{m}$ and $\rho_{\max}=5~\mathrm{m}$. The raw
target $\tilde{\rho}^{*}_t$ was filtered with a 0.1-s smoothing time to obtain
the applied target $\rho^{*}_t$. This absolute mapping allowed participants to
select formation size directly rather than controlling an expansion or
contraction rate.

\emph{FPV yaw control:}
Head motion controlled yaw rate rather than absolute view orientation. With
the neutral headset yaw denoted by $\psi_{0,t}$, the wrapped yaw deviation was
\[
\Delta\psi_t=
\operatorname{wrap}_{[-180^{\circ},180^{\circ}]}
\left(\psi_t-\psi_{0,t}\right).
\]
The resulting yaw-rate command was
\begin{equation}
\dot{\psi}^{\,\mathrm{cmd}}_t
=\omega_{\max}
\mathcal{D}\!\left(
\Delta\psi_t;\psi_{\mathrm{L}},0,\psi_{\mathrm{R}},10^{\circ}
\right),
\qquad
\omega_{\max}=48^{\circ}\mathrm{s}^{-1}.
\label{eq:yaw_rate}
\end{equation}
Returning the head to its neutral orientation therefore stopped view rotation
without forcing the view back to its initial heading. No additional yaw
smoothing was applied. To compensate for gradual postural drift,
$\psi_{0,t}$ was adapted at $0.5~\mathrm{s}^{-1}$ when head angular velocity
remained below $5^{\circ}\mathrm{s}^{-1}$ for at least 1.5~s and
$1^{\circ}\leq|\Delta\psi_t|\leq8^{\circ}$.


\subsection{Participant-Specific Calibration}
\label{sec:calibration}

\begin{figure}[!t]
\centering
\includegraphics[width=0.55\linewidth]{Figures/calib_coord.png}
\caption{Quantities measured during participant-specific calibration.
Controller positions are expressed in the Unity world frame
$\mathcal{F}_{W}$; their separation $s$ and mean vertical position $h$
provide the arm-spread and arm-height inputs. Head-yaw deviation
$\Delta\psi$ is measured about the headset frame $\mathcal{F}_{H}$, while
the chest-mounted IMU frame $\mathcal{F}_{C}$ provides the
forward--backward and lateral torso inclinations $\theta$ and $\phi$.}
\label{fig:participant_calibration}
\end{figure}

The interface used a guided 13-pose procedure to determine each participant's
neutral positions and comfortable motion limits. For each pose, an in-headset
cue provided 4~s to assume the requested posture, followed by a 3-s hold
countdown; the corresponding sensor value was recorded at the end of the hold.
Fig.~\ref{fig:participant_calibration} defines the angle conventions. The calibration sequence comprised
(i)~a neutral torso posture, which established the chest-IMU orientation
offset, followed by maximum-forward, maximum-backward, maximum-left, and
maximum-right inclinations, which defined the sagittal limits
$(\theta^{-},\theta^{+})$ and lateral limits
$(\phi^{-},\phi^{+})$ (five poses);
(ii)~a neutral head orientation, which established $\psi_0$, followed by
maximum-left and maximum-right turns, which defined
$(\psi_{\mathrm{L}},\psi_{\mathrm{R}})$ (three poses);
(iii)~minimum and maximum controller separations, which defined
$(s_{\min},s_{\max})$ (two poses); and
(iv)~minimum, neutral, and maximum controller heights, which defined
$(h_{\min},h_0,h_{\max})$ (three poses).
These parameters specified the normalizations in
Eq.~\eqref{eq:body_rate_commands}, the yaw normalization preceding
Eq.~\eqref{eq:yaw_rate}, and the absolute spread mapping in
Eq.~\eqref{eq:spread_mapping}. They were stored in a participant-specific
profile and used in the subsequent body-motion trials; no regression or
additional curve fitting was performed.

Table~\ref{tab:calibration_summary} summarizes the resulting calibrated input
spans. The observed between-participant variability supports the use of
participant-specific normalization rather than fixed population-wide bounds.

\begin{table}[t]
\centering
\caption{Distribution of participant-specific calibrated input spans.}
\label{tab:calibration_summary}
\footnotesize
\setlength{\tabcolsep}{3.5pt}
\begin{tabular}{@{}lccc@{}}
\toprule
Calibration measure & Mean $\pm$ SD & Range \\
\midrule
Sagittal torso ($^\circ$)
     & $56.4 \pm 18.4$  & $34.2$--$98.1$  \\
Lateral torso ($^\circ$)
     & $52.0 \pm 18.0$  & $26.2$--$81.2$  \\
Head yaw ($^\circ$)
    & $117.5 \pm 24.7$ & $65.6$--$138.0$ \\
Controller separation (m)
    & $1.17 \pm 0.31$  & $0.77$--$1.51$  \\
Controller height (m)
    & $0.77 \pm 0.22$  & $0.43$--$1.22$  \\
\bottomrule
\end{tabular}

\vspace{1mm}
\parbox{\columnwidth}{\scriptsize
Torso and head-yaw spans are the sums of the two directional excursion
magnitudes. Controller spans are the differences between the calibrated
maximum and minimum.}
\end{table}


\subsection{Joystick Baseline}
\begin{figure}[b]
\centering
\includegraphics[width=\columnwidth]{Figures/taranis_annotated.png}
\caption{Taranis transmitter and control-axis assignments used for the
joystick baseline. Straight arrows denote signed rate inputs; the curved arrow
denotes absolute control of target inter-agent spacing.}
\label{fig:taranis}
\end{figure}
\label{sec:joystick}

The baseline condition used a Taranis radio-control transmitter to command the
same five swarm DOF as the body-motion interface. Figure~\ref{fig:taranis}
summarizes the control-axis assignments.

Each of the four signed stick axes was normalized to $[-1,1]$ after applying
a symmetric deadzone of 0.12 and linearly rescaling the remaining range.
Translation and yaw commands used the same focal-drone reference frame,
translation gain $k_v=10~\mathrm{m\,s}^{-1}$, and maximum yaw rate
$\omega_{\max}=48^{\circ}\mathrm{s}^{-1}$ as the body-motion interface.
The right rotary knob bypassed the deadzone and mapped linearly from $[-1,1]$
to the same target-spacing range $\rho^{*}\in[1,5]~\mathrm{m}$.


% =====================================================================

\label{sec:experiment}

\section{Experimental Evaluation}
\label{sec:experiment}

\subsection{Participants and Study Design}
\label{sec:participants_design}

Fourteen participants (10 men and 4 women; age:
$22.4 \pm 4.0$ years, range 19--35) were recruited from the local university
community. Eligibility required no self-reported condition limiting the
necessary upper-body movements and vision compatible with the VR headset
without spectacles. Most participants had no drone-control experience,
although one was an experienced drone pilot; several reported prior video-game
experience. All participants were compensated for their time and provided
written informed consent.

A within-subject design was used to compare the body-motion and joystick
interfaces described in Section~\ref{sec:system}, thereby reducing the
influence of inter-participant variability. Both conditions were performed
while seated, as shown in Fig.~\ref{fig:setup}. Participants viewed the same
Unity environment through the same headset, and the simulated swarm,
focal-drone selection, command limits, and task were identical across
conditions. Thus, the conditions differed only in the physical input device
and its associated signal-processing pipeline. Condition order was
counterbalanced to limit order and learning effects, with seven participants
randomly assigned to each starting condition.

\begin{figure}[tb]
\centering
\includegraphics[width=\columnwidth]{Figures/set.png}
\caption{Experimental input conditions: (a) body-motion control using a Quest
headset, two Quest controllers, and a chest-mounted IMU; and (b) joystick
control using the same headset and a Taranis transmitter.}
\label{fig:setup}
\end{figure}

\subsection{Task}
\label{sec:task}


Both interfaces controlled the same 15-agent swarm described in
Section~\ref{sec:swarm_control}; all swarm-controller parameters were held
constant across conditions.
Three predefined course layouts were used: one practice course and two
recorded-trial courses. Each course
comprised seven wall gates and 316 collectible stars, and the layouts were
matched in nominal path length.

Each gate contained a single opening. Across the three courses, opening
dimensions ranged from $4.5\times3.75$~m to $15\times12.5$~m
(width $\times$ height), with opening centers displaced laterally and
vertically. Consecutive gates along straight segments were separated by
approximately 30~m, and each course turned by $90^{\circ}$ after the fifth
gate. This geometry required participants to coordinate planar translation,
altitude, view direction, and swarm contraction: lateral and vertical offsets
required three-dimensional steering, narrow openings required contraction,
and the final turn required a change in heading. The stars provided an
additional coverage objective along the route.

\begin{figure}[t]
\centering
\includegraphics[width=\columnwidth]{Figures/course_trials_annotated.png}
\caption{Recorded-trial course layouts: (a) trial~1 and (b) trial~2. Labels
G1--G7 indicate the gate sequence; cyan arrows show the traversal direction;
magenta and green planes mark the start and finish, respectively; and yellow
points represent collectible stars.}
\label{fig:course}
\end{figure}

A run began when the first swarm agent crossed the start plane and ended when
the first agent crossed the finish plane. Participants were instructed to
complete the course as quickly as possible while passing near the center of
each opening, collecting stars, and minimizing crashes and swarm
disconnections. A star was collected when any agent entered its trigger
volume. A crash was recorded when an agent collided with a wall or another agent, after which the affected agent was removed from the swarm. A disconnection
occurred when the swarm interaction graph split into multiple connected
components; agents outside the largest connected component were considered
disconnected.

\subsection{Procedure}
\label{sec:procedure}

Before the body-motion block, participants completed the participant-specific
calibration described in Section~\ref{sec:calibration}. Calibration was
performed once, in the same seated posture used during the trials, with no
recalibration between runs.

Each interface block began with a run on the practice course. Participants
could initially maneuver freely to familiarize themselves with the control
mapping, but all completed the practice course from start to finish. Practice
data were excluded from the analysis. Participants then completed two recorded
runs: trial~1 on course~1 and trial~2 on course~2. The course-to-trial
assignment was identical across interface conditions to provide equivalent
paired tasks, while the separate practice course prevented practice from
affecting the recorded courses.

After each interface block, participants completed the NASA Task Load Index
(NASA-TLX)~\cite{hart1988nasatlx} and System Usability Scale
(SUS)~\cite{brooke1996sus}. After both blocks, they completed a comparative
questionnaire concerning interface preference and relative motion sickness.

\subsection{Evaluation Metrics}

We defined the following metrics to compare traversal efficiency, delivered
commands, task outcomes, and subjective experience.

\emph{Traversal efficiency.}
Completion time $T$ was the interval between the start and finish events.
For the controlled-component centroid positions $\mathbf{x}_k$, we also
calculated travel distance
$L=\sum_k\lVert\mathbf{x}_{k+1}-\mathbf{x}_k\rVert_2$ and mean centroid speed
$\bar{v}=L/T$. Completion time was the primary task outcome; speed and distance
were supporting measures because they are mathematically related to $T$.

\emph{Delivered-command behavior.}
The five delivered commands (three translations, target spacing, and view-yaw
rate) were normalized by their implemented ranges to
$\mathbf{q}_k\in[0,1]^5$. With $\tau=t_M-t_1$,
$\Delta\mathbf{q}_k=\mathbf{q}_{k+1}-\mathbf{q}_k$, and
$\Delta t_k=t_{k+1}-t_k$, command variation was
\begin{equation}
E_{\mathrm{cmd}}=\frac{\tau}{5}\sum_{k=1}^{M-1}
\frac{\lVert\Delta\mathbf{q}_k\rVert_2^2}{\Delta t_k}.
\label{eq:command_metrics}
\end{equation}
Lower $E_{\mathrm{cmd}}$ indicates less variation in the commands supplied to
the swarm. Because the interfaces used different online filters, this metric
characterizes the complete delivered-command pipelines rather than the
intrinsic smoothness of the input modalities.

\emph{Task accuracy and swarm integrity.}
At the first forward crossing within each physical gate, centroid accuracy was
\begin{equation}
E_{\mathrm{gate}}=\frac{1}{7}\sum_{g=1}^{7}
\max\left(\frac{|h_g|}{W_g/2},\frac{|v_g|}{H_g/2}\right),
\label{eq:gate_error}
\end{equation}
where $h_g$ and $v_g$ are offsets from the opening center and $W_g$ and $H_g$
are its dimensions. We further measured unique-star yield, swarm survival
($1-N_{\mathrm{crash}}/15$), and the time-weighted fraction of surviving
agents in the largest connected component.

\emph{Subjective measures.}
NASA-TLX and SUS quantified workload and usability. The final questionnaire
recorded interface preference and relative motion sickness. We analyzed the
recorded weighted TLX total and its six 0--100 ratings. Because TLX importance
weights were elicited separately after each interface, the unweighted mean of
the six ratings (Raw TLX) was also computed as a scoring sensitivity.

\subsection{Data Processing and Statistical Analysis}
\label{sec:statistical_analysis}

All 56 recorded runs were retained. For each metric defined above, the 15-Hz
logs were restricted to the start--finish interval; boundary values were
interpolated, identical duplicate timestamps removed, and no smoothing or
resampling applied. A value was computed per run and the two trials were then
averaged within participant and interface.

The resulting paired values were compared using exact two-sided Wilcoxon
signed-rank tests, with Holm correction within the efficiency, command,
task-outcome, and global-subjective families. The six exploratory TLX
dimensions formed a separate family. Figure~\ref{fig:results_summary} reports
interface means, paired body-minus-joystick differences, bootstrap 95\%
confidence intervals (20,000 resamples), and adjusted $p$-values. Sensitivity
checks examined trial/course and condition-order effects, path resampling,
common filtering of both command streams, and Raw TLX scoring.

% =====================================================================
\section{Results}

All 14 participants contributed two valid runs per interface.
Figure~\ref{fig:results_summary} combines the statistical summary with
participant-level results for completion time and physical demand.

\emph{Traversal efficiency.}
Body-motion control reduced mean completion time from 115.3 to 92.6~s, a
19.7\% reduction ($\Delta=-22.7$~s; 95\% CI $[-32.1,-13.6]$;
$p_{\mathrm H}=.001$). Thirteen of the 14 participants were faster with body
motion (Fig.~\ref{fig:results_summary}(a)). Their trajectories were also 6.6\%
shorter ($\Delta=-21.2$~m; 95\% CI $[-34.4,-7.4]$;
$p_{\mathrm H}=.011$). Correspondingly, mean centroid speed increased from
2.90 to 3.28~m/s, a 13.2\% increase.

\emph{Delivered-command behavior.}
Delivered command variation was 88.7\% lower in the body-motion pipeline
($\Delta=-2971$; 95\% CI $[-3938,-2075]$;
$p_{\mathrm H}<.001$). Applying the same additional 0.143-s low-pass filter
to both delivered streams preserved a 79.8\% reduction. However, unfiltered
body commands were not logged, so this result describes the complete delivered
pipelines rather than an intrinsic difference between the two input signals.

\emph{Task outcomes.}
Neither normalized gate error nor collection yield showed a clear interface
difference. Swarm survival and largest-component retention remained near
ceiling with both interfaces: survival averaged 99.0\% with body motion and
98.6\% with the joystick, while component retention averaged 0.996 and 1.000,
respectively. These results do not establish equivalent task performance.

\emph{Subjective outcomes.}
Weighted NASA-TLX and SUS showed no clear interface differences; Raw TLX gave
the same conclusion. Physical demand, however, was higher with body motion for
all 14 participants, increasing from 23.0 to 63.7 points
($\Delta=40.7$; 95\% CI $[30.7,50.4]$; $p_{\mathrm H}<.001$;
Fig.~\ref{fig:results_summary}(b)). Eight participants preferred body motion,
five preferred the joystick, and one reported no preference. Six reported
greater motion sickness with body motion, four with the joystick, and four
reported no difference.

\begin{strip}
\centering

\begin{minipage}[t]{0.55\textwidth}
\vspace{0pt}
\centering
\scriptsize
\setlength{\tabcolsep}{2pt}
\renewcommand{\arraystretch}{1.04}

\begin{tabular}{@{}lcc@{}}
\toprule
Outcome &
Body/joystick &
$\Delta$ [95\% CI]; $p_{\mathrm H}$ \\
\midrule

\multicolumn{3}{@{}l}{\emph{Traversal efficiency}} \\
Completion time (s) &
$92.6/115.3$ &
$-22.7\,[-32.1,-13.6];\ .001$ \\

Path length (m) &
$298.9/320.1$ &
$-21.2\,[-34.4,-7.4];\ .011$ \\

\addlinespace[2pt]
\multicolumn{3}{@{}l}{\emph{Delivered-command behavior}} \\
Command variation$^{*}$ &
$380/3351$ &
$-2971\,[-3938,-2075];\ <.001$ \\

\addlinespace[2pt]
\multicolumn{3}{@{}l}{\emph{Task outcomes}} \\
Normalized gate error &
$0.236/0.228$ &
$0.008\,[-0.023,0.039];\ 1.000$ \\

Collection yield (\%) &
$25.6/24.8$ &
$0.8\,[-2.1,3.6]$ pp; $1.000$ \\

\addlinespace[2pt]
\multicolumn{3}{@{}l}{\emph{Subjective outcomes}} \\
Weighted NASA-TLX &
$59.1/51.9$ &
$7.2\,[-5.5,20.2];\ .482$ \\

TLX physical demand &
$63.7/23.0$ &
$40.7\,[30.7,50.4];\ <.001$ \\

SUS &
$60.7/59.1$ &
$1.6\,[-12.1,14.1];\ .491$ \\

\bottomrule
\end{tabular}
\end{minipage}
\hfill
\begin{minipage}[t]{0.42\textwidth}
\vspace{0pt}
\centering
\includegraphics[width=\linewidth]{Figures/fig_key_results.pdf}
\end{minipage}

\captionof{figure}{Summary of paired interface outcomes ($n=14$). Left:
interface means (body/joystick), paired body-minus-joystick differences with
bootstrap 95\% CIs, and Holm-adjusted $p_{\mathrm H}$ values. Right:
participant-level completion time (a) and TLX physical demand (b) after
averaging the two trials per interface; circles denote participants, lines
connect conditions, and diamonds show group means. Confidence intervals
containing zero do not establish equivalence. $^{*}$Command variation was
measured at the delivered-command output and includes interface-specific
filtering.}
\label{fig:results_summary}
\end{strip}

\FloatBarrier


% =====================================================================
\section{Discussion}
% --- 4. Discussion ---

Within the tested simulated FPV task, the body-motion interface improved
traversal efficiency relative to the conventional transmitter. Participants
completed the courses 19.7\% faster while following shorter centroid
trajectories and attaining a higher mean centroid speed. No interface
difference was detected in gate centering, collection yield, swarm survival,
or largest-component retention. The results therefore support an efficiency
advantage for the complete implemented body-motion interface under the
evaluated conditions, rather than a general advantage across all outcomes.

The reduction in completion time reflected both faster motion and more direct
trajectories. A plausible explanation is that distributing commands across the
head, torso, and arms allowed several swarm DOF to be modified concurrently
while maintaining a spatially congruent mapping in the rotating focal-drone
frame. By contrast, adjusting formation spread with the transmitter required
participants to leave the forward-velocity stick to manipulate the spread
knob, which may have serialized commands and introduced pauses near narrow
gates. The higher simultaneous-change index under body-motion control is
consistent with the concept of integral multidimensional input
\cite{jacob1994integrality,zhai1998quantifying}. However, this index measures
temporal concurrence rather than coordination quality, and the present
experiment does not establish that concurrent commands caused the performance
improvement. Similarly, the lower delivered-command variation characterizes
the complete interface pipelines: the body-motion interface applied online
filtering, whereas the transmitter relied primarily on a deadzone. Matched
filtering or a filter-ablation study would be required to isolate the effect of
the input modality.

The efficiency improvement was not accompanied by a detectable degradation in
the measured task-quality or swarm-integrity outcomes. These null results do
not demonstrate equivalence or noninferiority. In particular, swarm survival
and largest-component retention were close to their upper bounds in both
conditions, limiting sensitivity to modest differences. Consequently, the
findings should not be interpreted as evidence of equivalent safety or
robustness beyond the tested simulation.

The subjective outcomes reveal a performance--effort trade-off: body-motion
control improved objective efficiency but was substantially more physically
demanding, while total NASA-TLX and SUS did not differ detectably between
interfaces. Thus, its performance benefit came at a physical cost without a
clear overall workload or usability advantage.

Overall, these findings extend upper-body interfaces previously evaluated for
individual-drone control
\cite{rognon2018flyjacket,miehlbradt2018datadriven,macchini2024datadriven, tsykunov2019swarmtouch,kratky2025gesture}
and complement hand- or body-based swarm interfaces evaluated primarily from
external viewpoints
\cite{macchini2021personalized,tsykunov2019swarmtouch}. They also complement
previous FPV swarm research using a conventional transmitter
\cite{ben_focaldrone}. By holding the FPV strategy and swarm controller
constant, the present comparison indicates that distributing collective
control dimensions across body segments is a promising approach when
translation, altitude, view direction, and formation size must be adjusted
during the same maneuver. Direct quantitative comparison with previous studies
is nevertheless inappropriate because their tasks, viewpoints, swarm
configurations, and outcome measures differ.

The physical-demand result identifies an important direction for interface
refinement. FlyJacket incorporated passive arm support to reduce fatigue during
upper-body drone control \cite{rognon2018flyjacket}. A comparable mechanism
could be investigated for the present interface. Participants also reported
that torso leaning could unintentionally lower their hands and therefore
affect the vertical command. Passive support or software compensation that
separates torso-induced hand displacement from intentional arm-height changes
could reduce this command coupling. Informal participant feedback further
indicated limited awareness of the swarm's spatial extent when selecting a
formation spread for passing through openings. Vibrotactile feedback could
communicate the current spread or available clearance, as demonstrated for
swarm-state and formation guidance in SwarmTouch
\cite{tsykunov2019swarmtouch}. These observations were not obtained through a
formal qualitative analysis and should therefore be treated as design
hypotheses for future evaluation.

Finally, the findings are bounded by the relatively small and predominantly
young university sample, two recorded courses with a fixed course-to-trial
order, one 15-agent swarm, one rear-FPV strategy, seated operation, and brief
exposure to each interface. The simulation also omitted real world noise (wind, communication
disturbances, consequences of physical collisions, ...). Future studies should examine learning and fatigue over longer periods, randomize and diversify the courses, evaluate different swarm
sizes and viewpoint strategies, compare interfaces under matched signal
processing, and validate the approach using physical aerial swarms.

% =====================================================================
\section{Conclusion}

This work presented a participant-calibrated upper-body interface for
continuous five-DOF teleoperation of an aerial swarm under first-person view.
In a within-subject comparison with a conventional transmitter, the complete
body-motion interface enabled participants to finish the simulated courses
nearly 20\% faster while travelling shorter distances and attaining higher mean
speed. It also produced lower delivered-command variation and more
simultaneous multi-DOF command changes. No interface differences were detected
in gate centering, collection, swarm survival, connectivity, workload, or
usability, while preference and motion-sickness responses were mixed. These
findings support upper-body motion as a promising command modality for
multi-DOF FPV swarm interaction within the tested simulation. Evaluation under
matched signal processing, longer exposure, varied swarm configurations, and
physical flight is required before broader deployment conclusions can be
drawn.


% =====================================================================
\bibliographystyle{IEEEtran}
\bibliography{references}

\end{document}
